\documentclass[10pt,a4paper]{amsart}
\usepackage[margin=19mm]{geometry}
\usepackage{amsmath,amssymb,mathtools}
\usepackage[scr=rsfs,cal=euler]{mathalfa}
\usepackage{xfrac}
\usepackage[unicode,hidelinks,bookmarks=false]{hyperref}
\newcommand{\ie}{i.e.\ }
\newcommand{\cf}{cf.\ }
\newcommand{\resp}{respectively\ }
\newcommand{\Subsection}{\subsection}
\providecommand{\nobreakdash}{\nobreak\hbox{-}}
\setlength{\emergencystretch}{2em}
\multlinegap0pt
\usepackage{bm}
\usepackage{bbm}
\usepackage{dsfont}
\usepackage{xspace}
\usepackage{cancel}
\usepackage{mathtools}
\usepackage{amsthm}
\makeatletter
\@ifundefined{theorem}{\newtheorem{theorem}{Theorem}}{}
\@ifundefined{cor}{\newtheorem{cor}[theorem]{Corollary}}{}
\@ifundefined{defin}{\newtheorem{defin}[theorem]{Definition}}{}
\@ifundefined{prop}{\newtheorem{prop}[theorem]{Proposition}}{}
\makeatother
\def\cI{\mathcal I}
\def\cJ{\mathcal J}
\def\cP{\mathcal P}
\def\cT{\mathcal T}
\def\cU{\mathcal U}
\title[On line-parallelisms of PG(3, q)]{On line-parallelisms of PG(3, q)}
\author{Francesco Pavese, Paolo Santonastaso}
\date{Source: arXiv:2506.16271v1 (CC BY 4.0)}
\begin{document}
\maketitle

\noindent\emph{Source and licence.} On line-parallelisms of PG(3, q). Version arXiv:2506.16271v1 (CC BY 4.0), distributed under the Creative Commons Attribution 4.0 International licence, as recorded in the arXiv licence metadata for this exact version and shown on the abstract page https://arxiv.org/abs/2506.16271v1. Full rights notice: licenses/arxiv-2506.16271-CC-BY-4.0.txt.\par\smallskip

\noindent\textit{Editorial scope.} Counted unit: the Theorem whose source label is \texttt{None} in the article (source0.tex lines 1473--1480, source label \texttt{None}), delivered together with the article's own complete original proof of it (source0.tex lines 1481--1533, one \texttt{proof} environment of 53 lines). No mathematics is written, repaired or re-derived: statement and proof are the article's own text line for line. Required context retained uncounted: goodset (lines 883--892); prop:allgoodsetseven (lines 1258--1263); cor:allgoodsetsodd (lines 1341--1351); automorphisms (lines 1397--1399); distinctgoodsets (lines 1449--1461). Every definition, display and label of those lines is retained unchanged. Omitted from this package and not redistributed: 1--882, 893--1257, 1264--1340, 1352--1396, 1400--1448, 1462--1472, 1534--1659. Nothing inside a retained excerpt is omitted or altered: every retained body is delivered exactly as the article's lines are written, so no omission receipt of any kind accompanies this package. Citations: the retained excerpts carry no citation command at all, so no bibliography entry of the article is transcribed and no reference is redistributed. Reference adaptation: none was needed and none was made; every reference inside the delivered unit resolves inside it, so no unresolved reference or citation is produced. Printed slips kept literal: no printed word, symbol, hypothesis or number of the counted statement or of its proof was altered. Layout and class: the article's own class is \texttt{article} with packages including amsmath, amsthm, amssymb, color, hyperref, pxfonts, mathrsfs, titlesec; the wrapper is the lane's pinned \texttt{amsart} editorial wrapper at 10pt on A4, which restates the theorem-like environments the retained text needs and adds the article's own macro definitions that the retained text needs (\texttt{\textbackslash cI}, \texttt{\textbackslash cJ}, \texttt{\textbackslash cP}, \texttt{\textbackslash cT}, \texttt{\textbackslash cU}), with extra wrapper packages (bm, bbm, dsfont, xspace, cancel, mathtools). The delivered numbering is the unit's own. Assets: no retained span contains a figure environment, a graphics-inclusion command, a PDF-page inclusion, a table or a TikZ picture; no image file is copied into this package, the package declares no asset dependency and no image file is redistributed with it. Licence: version arXiv:2506.16271v1 of this article is distributed under the Creative Commons Attribution 4.0 International licence, as recorded in the arXiv licence metadata for this exact version and shown on the abstract page https://arxiv.org/abs/2506.16271v1. A full rights notice is delivered at licenses/arxiv-2506.16271-CC-BY-4.0.txt. Characters: every retained excerpt is pure ASCII, so the delivered engine, XeLaTeX with its default Latin Modern text font, reports no missing character.
\begin{defin}\label{goodset}
    A set
\begin{align*}
& \cP = \left\{\left(P_{\alpha_{1} u_1}, \pi_{\alpha_{1} v_1}\right), \dots, \left(P_{\alpha_{{q+1}} u_{q+1}}, \pi_{\alpha_{{q+1}} v_{q+1}}\right)\right\}, & \alpha_{i} \in \cI, u_i, v_i \in \cU, 
\end{align*}
consisting of $q+1$ pairs, is said to be a {\em good set} if $\forall \alpha_{i}, \alpha_{j} \in \cI$, $\forall u_i, u_j, v_i, v_j \in \cU$, with $(\alpha_{i}, u_{i}, v_i) \ne (\alpha_{j}, u_j, v_j)$, it satisfies the following properties:
\begin{align}
\left(P_{\alpha_{i} u_i}, \pi_{\alpha_{i} v_i}\right), \left(P_{\alpha_{j} u_{j}}, \pi_{\alpha_{j} v_j}\right) \in \cP & \implies u_i v_j - v_i u_j \ne 0, \label{property1} \\
\left(P_{\alpha_{i} u_i}, \pi_{\alpha_{i} v_i}\right), \left(P_{\alpha_{j} u_{j}}, \pi_{\alpha_{j} v_{j}}\right) \in \cP & \implies \alpha_{i} u_i \left(\alpha_{j} v_{j}\right)^q - \left(\alpha_{i} v_i\right)^q \alpha_{j} u_{j} \ne 0. \label{property2}\end{align}
\end{defin}

\begin{prop}\label{prop:allgoodsetseven}
    Assume that $q$ is even. The number of good sets equals  
    \begin{align*}
    \lvert \cI \rvert^{q+1} (q+1)!=\left( \frac{q-1}{2} \right)^{q+1} \cdot (q+1)!
    \end{align*}
\end{prop}

\begin{cor}\label{cor:allgoodsetsodd}
Assume that $q$ is odd. The number of good sets 
\begin{align*}
& \cP = \left\{\left(P_{\alpha_{1} u_1}, \pi_{\alpha_{1} v_1}\right), \dots, \left(P_{\alpha_{{q+1}} u_{q+1}}, \pi_{\alpha_{{q+1}} v_{q+1}}\right)\right\},  
\end{align*}
with $\alpha_i^{q+1} \neq -1$, $i = 1, \dots, q+1$, equals  
    \begin{align*}
    & \left(\frac{(q-5)(q-1)}{16}\right)^{\frac{q+1}{2}} \prod_{i=0}^{\frac{q-1}{2}} \left(q+1-2i \right)^2 & \mbox{ if }q \equiv 1 \pmod{4}, \\ 
    & \left(\frac{q-3}{4}\right)^{q+1} \prod_{i=0}^{\frac{q-1}{2}} \left(q+1-2i \right)^2 & \mbox{ if } q \equiv 3 \pmod{4}. 
    \end{align*}
\end{cor}

\begin{prop}\label{automorphisms}
Let $\Pi_{\cP}$, $\Pi_{\cP'}$ line-parallelisms of $\Sigma_{\eta}$, where $\cP$, $\cP'$ are good sets. If there exists $ \psi \in \cJ$ such that $\Pi_{\cP}^\psi = \Pi_{\cP'}$, then $\psi \in \Gamma_{r_{U_1}}$. 
\end{prop}

\begin{prop}\label{distinctgoodsets}
%    Assume that $\alpha_2^{q+1} \ne -1$, in the case when $q$ is odd. Then 
%\begin{align*}
Let \begin{align*}
& \cP = \left\{\left(P_{\alpha_{1} u_1}, \pi_{\alpha_{1} v_1}\right), \dots, \left(P_{\alpha_{{q+1}} u_{q+1}}, \pi_{\alpha_{{q+1}} v_{q+1}}\right)\right\}, & \alpha_{i} \in \cI, u_i, v_i \in \cU, 
\end{align*}
and 
\begin{align*}
& \cP' = \left\{\left(P_{\alpha_{1}' u_1'}, \pi_{\alpha_{1}' v_1'}\right), \dots, \left(P_{\alpha_{q+1}' u'_{q+1}}, \pi_{\alpha_{q+1}' v'_{q+1}}\right)\right\}, & \alpha_{i}' \in \cI, u_i', v_i' \in \cU, 
\end{align*}
be good sets. Assume that $\alpha_i^{q+1}\neq -1$ and $\alpha_i'^{q+1}\neq -1$, $i = 1, \dots, q+1$. If $\cP \ne \cP'$, then $\Pi_{\cP} \ne \Pi_{\cP'}$.
%\end{align*} 
\end{prop}

\begin{theorem}
The number of mutually not equivalent line-parallelisms of $\Sigma_{\eta}$ contained in $\cT$ is at least
    \begin{align*}
      & \left(\frac{q-1}{2}\right)^{q+1} \frac{(q-2)!}{2hq(q+1)} & \mbox{ if } q \mbox{ is even, } \\
      & \left(\frac{(q-5)(q-1)}{16}\right)^{\frac{q+1}{2}} \frac{(q^2-1)^{\frac{q-1}{2}}}{2hq^2(q+1)} & \mbox{ if }q \equiv 1 \pmod{4}, \\ 
      &  \left(\frac{q-3}{4}\right)^{q+1} \frac{(q^2-1)^{\frac{q-1}{2}}}{2hq^2(q+1)} & \mbox{ if } q \equiv 3 \pmod{4}.  
    \end{align*}
\end{theorem}

\begin{proof}
From Proposition~\ref{automorphisms}, if $\Pi_{\cP}, \Pi_{\cP'} \in \cT$ and $\Pi_{\cP}^\psi = \Pi_{\cP'}$, for some $\psi \in \cJ$, then $\psi \in \Gamma_{r_{U_1}}$. Hence
\begin{equation} \label{eq:boundstabiliser}
|\Pi_{\cP}^{\cJ} \cap \cT| \le |\Gamma_{r_{U_1}}| = 2hq^2(q^2-1)(q+1)
\end{equation}

Assume first that $q$ is even. By the definition of good set (see Definition \ref{goodset}), we know that each good set
\begin{align*}
\cP = \left\{\left(P_{\alpha_{1} u_1}, \pi_{\alpha_{1} v_1}\right), \dots, \left(P_{\alpha_{q+1} u_{q+1}}, \pi_{\alpha_{q+1} v_{q+1}}\right)\right\}
\end{align*}
is such that $\alpha_i^{q+1} \neq 1$, $i = 1, \dots, q+1$. Therefore, by Proposition \ref{distinctgoodsets}, the number of distinct parallelisms in $\cT$ equals the number of distinct good sets. By Proposition \ref{prop:allgoodsetseven}, this number is given by:
\[
|\cT| = |\cI|^{q+1} \cdot (q+1)! = \left( \frac{q-1}{2} \right)^{q+1} \cdot (q+1)!.
\]
Now, considering \eqref{eq:boundstabiliser}, we obtain that in $\cT$ there are at least
\[
\begin{array}{rl}
\displaystyle \frac{|\cT|}{|\Pi_{\cP}^{\cJ} \cap \cT|} &\displaystyle \ge \left( \frac{q-1}{2} \right)^{q+1} \cdot \frac{(q+1)!}{2hq^2(q^2 - 1)(q+1)} \\[10pt]
&\displaystyle = \left( \frac{q-1}{2} \right)^{q+1} \cdot \frac{(q-2)!}{2hq(q + 1)}
\end{array}
\]
non-equivalent line-parallelisms of $\Sigma_{\eta}$.


Now, assume that $q$ is odd. The size of $\cT$ is clearly greater than or equal to the number of good sets
\begin{align*}
\cP = \left\{\left(P_{\alpha_{1} u_1}, \pi_{\alpha_{1} v_1}\right), \dots, \left(P_{\alpha_{q+1} u_{q+1}}, \pi_{\alpha_{q+1} v_{q+1}}\right)\right\},
\end{align*}
having the property that $\alpha_i^{q+1} \neq -1$, $i = 1, \dots, q+1$. Therefore, by Corollary \ref{cor:allgoodsetsodd} and taking into account Proposition \ref{distinctgoodsets}, we obtain
\[
|\cT| \geq 
    \begin{cases} 
    \left(\frac{(q-5)(q-1)}{16}\right)^{\frac{q+1}{2}} \prod\limits_{i=0}^{\frac{q-1}{2}} \left(q+1-2i \right)^2 & \mbox{ if }q \equiv 1 \pmod{4}, \\ 
    \left(\frac{q-3}{4}\right)^{q+1} \prod\limits_{i=0}^{\frac{q-1}{2}} \left(q+1-2i \right)^2 & \mbox{ if } q \equiv 3 \pmod{4}. 
    \end{cases}
\]
So, when $q$ is odd, in $\cT$ there are at least 
\begin{align*}
    \frac{|\cT|}{|\Pi_{\cP}^{\cJ} \cap \cT|} \ge 
   \begin{cases} 
   \left(\frac{(q-5)(q-1)}{16}\right)^{\frac{q+1}{2}} \frac{(q^2-1)^{\frac{q-1}{2}}}{2hq^2(q+1)} & \mbox{ if }q \equiv 1 \pmod{4} \\ 
   \left(\frac{q-3}{4}\right)^{q+1} \frac{(q^2-1)^{\frac{q-1}{2}}}{2hq^2(q+1)} & \mbox{ if } q \equiv 3 \pmod{4}
   \end{cases}
\end{align*}
not equivalent line-parallelisms of $\Sigma_{\eta}$. Indeed,
\begin{align*}
    \frac{\prod_{i = 0}^{\frac{q-1}{2}} \left(q+1 - 2 i \right)^2}{2hq^2(q^2-1)(q+1)} > \frac{(q^2-1)^{\frac{q-1}{2}}}{2hq^2(q+1)},
\end{align*}
since
\begin{align*}
    \prod_{i = 0}^{\frac{q-1}{2}} \left(q+1 - 2i \right)^2 > \prod_{i = 0}^{\frac{q-1}{2}} \left((q+1)^2 - 2 (q+1) \right) = \prod_{i = 0}^{\frac{q-1}{2}} \left(q^2-1 \right) = \left(q^2-1\right)^\frac{q+1}{2}.
\end{align*}
\end{proof}

\end{document}
